\documentclass[12pt]{article}
\usepackage[utf8]{inputenc}
\usepackage[T1]{fontenc}
\usepackage{hyperref}
\title{Mori Bawah Language Resources for K-12 Curriculum Development}
\author{Compiled by Claude Sonnet 5.5 (claude-sonnet-5-5) for Andr\'as Kornai}
\date{}
\begin{document}
\maketitle

\section{Introduction}
Mori Bawah (ISO 639-3 xmz) is an Austronesian language of the Bungku-Tolaki group, spoken in Morowali Regency, Central Sulawesi, Indonesia. Its main dialects are Tinompo, Tiu, Moiki, Watu and Karunsi'e. Speaker estimates vary by source: about 28,000 (a 1988 figure shared with Mori Atas), 14,000, and 22,000. The search found little material for school use.

\section{Primary Educational Resources}
I found no primers, textbooks or school programmes for Mori Bawah.

\section{Linguistic and Reference Materials}
\begin{itemize}
\item Wikipedia, Mori Bawah language: \url{https://en.wikipedia.org/wiki/Mori_Bawah_language}
\item Wikipedia, Mori language: \url{https://en.wikipedia.org/wiki/Mori_language}
\item Omniglot, Mori Bawah. It says the language is written in Latin script and that a grammar, some texts and a partial Bible translation exist: \url{https://www.omniglot.com/writing/moribawah.htm}
\item Austronesian Basic Vocabulary Database, Mori Bawah (Watu) wordlist: \url{https://abvd.eva.mpg.de/austronesian/language.php?id=902}
\item Verbix language page: \url{https://docs.verbix.com/Languages/MoriBawah}
\end{itemize}

\section{Religious and Community Resources}
\begin{itemize}
\item Joshua Project, Mori Bawah people group profile: \url{https://www.joshuaproject.net/people_groups/13814/ID}
\item Global Recordings Network, Mori Bawah audio Bible stories and gospel material: \url{https://globalrecordings.net/en/language/708}
\end{itemize}

\section{Digital and Audio Resources}
The Global Recordings Network page above is the only audio source found.

\section{Curriculum Development Context}
The language is used by adults and is not known to be taught in schools. The sources I found are reference-level only. Any curriculum would need fieldwork with the community, and Tinompo is the dialect with the most prestige. I did not verify the existence or location of the grammar and texts that Omniglot mentions.
\end{document}